%!TEX root = ../thesis.tex
% ===================================================================
% Chapter 4: Implementation (replaces the stubs in chapters/implementation.tex)
% Uses the same tikzset as design_chapter.tex. If design.tex already defines
% it (it does, as of 11 Oct), delete the \tikzset block below to avoid a
% duplicate-definition warning; otherwise keep it, since this chapter draws
% a set of block diagrams (fig:providers, fig:protocol) plus a flowchart per
% major function (dec = decision diamond, proc = process box, term = start/end).
% ===================================================================
\mychapter{Implementation}{Implementation}{}
\label{chap:implementation}

\tikzset{
  blk/.style={draw, rounded corners=2pt, align=center, font=\footnotesize,
              minimum width=2.6cm, minimum height=0.95cm, inner sep=3pt},
  mdl/.style={blk, fill=black!8},
  hw/.style={blk, dashed},
  st/.style={draw, rounded corners=2pt, align=center, font=\footnotesize,
             minimum width=2.7cm, minimum height=0.9cm},
  arr/.style={-{Stealth[length=2mm]}, thick},
  lbl/.style={font=\scriptsize, fill=white, inner sep=1pt},
  proc/.style={draw, rounded corners=2pt, align=center, font=\scriptsize,
               minimum width=2.9cm, minimum height=0.6cm, inner sep=2pt},
  dec/.style={draw, diamond, aspect=2, align=center, font=\scriptsize,
              inner sep=2pt, minimum width=2.6cm, minimum height=1.05cm},
  term/.style={draw, rounded corners=9pt, align=center, font=\scriptsize,
               minimum width=2.4cm, minimum height=0.55cm, inner sep=2pt},
  fa/.style={-{Stealth[length=1.8mm]}, thin},
  fl/.style={font=\tiny, fill=white, inner sep=1pt}
}

Chapter~\ref{chap:design} said what each part of the pipeline does and why.
This chapter says how each part is actually built: the libraries, the process
layout, the message formats, and the places where a design decision turned
into a specific piece of code. I keep the numbers here that
Chapter~\ref{chap:design} left out on purpose, such as port numbers, file
names and exact thresholds, since those belong to how the system is built
rather than why it is shaped the way it is.

Every class in this chapter is described in a short paragraph, since a class
is a data shape, not a behaviour, and reads better as prose than as a
diagram. Every function that actually branches is instead described with a
small flowchart: a diamond is a decision, a rounded rectangle is a step, and
the text around the figure covers what the shapes cannot, such as exact
constants and file names.

%% ===================================================================
\mysection{Speech input and output}{Speech Input and Output}
\label{sec:speech}

\mysubsection{Speech-to-text}{Speech-to-text}
\label{sec:stt}

Transcription runs through \texttt{faster-whisper}, a CTranslate2
reimplementation of Whisper, using the \texttt{digiphyte/\allowbreak
fluister-turbo} model. I picked that model over stock Whisper because it is
trained to include South African English and Afrikaans accents, which
matters for a project the supervisor will test in person. The model runs on
the CPU with \texttt{compute\_type="int8"} and 8 threads, kept off the GPU on
purpose since the planner and the vision model already share the 12\,GB of
video memory. \texttt{transcribe\_audio} takes either raw bytes (an encoded
clip, decoded through FFmpeg) or a NumPy array of raw float32 samples, since
the capture path that reads straight from a microphone has no container
format to decode, and is called once per utterance rather than as a stream,
since Whisper pads every clip to 30\,s internally regardless.

\begin{figure}[htbp]
  \centering
  \begin{tikzpicture}
    \node[term] (a) at (0,0)    {audio in};
    \node[proc] (b) at (0,-0.9) {decode to float32\\(bytes: via FFmpeg)};
    \node[proc] (c) at (0,-1.8) {\texttt{faster-whisper}\\transcribe};
    \node[proc] (d) at (0,-2.7) {lowercase,\\strip punctuation};
    \node[dec]  (e) at (0,-4.1) {matches a fixed\\hallucination\\phrase?};
    \node[term] (f) at (-2.7,-5.6) {return \texttt{""}};
    \node[term] (g) at (2.7,-5.6)  {return transcript};
    \draw[fa] (a)--(b); \draw[fa] (b)--(c); \draw[fa] (c)--(d); \draw[fa] (d)--(e);
    \draw[fa] (e.south west) -- node[fl,above]{yes} (f.north);
    \draw[fa] (e.south east) -- node[fl,above]{no} (g.north);
  \end{tikzpicture}
  \caption{\texttt{transcribe\_audio}. The hallucination check exists because \texttt{fluister-turbo} occasionally transcribes silence or background noise as one of a handful of fixed captioning artefacts, most often ``thank you so much for watching''; I found this by reading failed-command logs where the ambiguity check was asked to resolve a command that was never actually said.}
  \label{fig:flow-stt}
\end{figure}

Returning an empty string rather than the artefact phrase means the rest of
the pipeline treats a hallucinated caption exactly like a transcription that
genuinely failed, rather than as a real command to act on.

\mysubsection{Speech output: Piper}{Speech Output}
\label{ssec:tts}

Spoken output goes through Piper, a local neural TTS engine, using the
\texttt{en\_US-lessac-medium} voice. \texttt{tts.py} mirrors \texttt{stt.py}'s
cached-model pattern: \texttt{get\_voice} loads the \texttt{PiperVoice} once
from \texttt{PIPER\_VOICE\_MODEL} and keeps it in a module-level variable.
Piper's voices are not fetched automatically the way \texttt{faster-whisper}'s
model is; I had to run \texttt{python -m piper.\allowbreak download\_voices
en\_US-\allowbreak lessac-medium} once by hand and point the environment
variable at the resulting \texttt{.onnx} file, with its paired
\texttt{.onnx.json} config sitting alongside it. There is no Afrikaans voice
in the \texttt{rhasspy/\allowbreak piper-voices} set, which I checked rather
than assumed; this turned out not to matter, since every \texttt{\_speak}
call site speaks a hard-coded English template string regardless of the
session's detected language, so only the speech-to-text side needs to branch
on language at all.

\begin{figure}[htbp]
  \centering
  \begin{tikzpicture}
    \node[term] (a) at (0,0)    {\texttt{\_speak(text)}};
    \node[proc] (b) at (0,-0.9) {\texttt{asyncio.to\_thread}\\(\texttt{synthesize\_speech})};
    \node[dec]  (c) at (0,-2.3) {synthesis\\throws?};
    \node[term] (d) at (-2.7,-3.8) {event: text only\\(log a warning)};
    \node[term] (e) at (2.7,-3.8)  {event: text +\\base64 audio};
    \draw[fa] (a)--(b); \draw[fa] (b)--(c);
    \draw[fa] (c.south west) -- node[fl,above]{yes} (d.north);
    \draw[fa] (c.south east) -- node[fl,above]{no} (e.north);
  \end{tikzpicture}
  \caption{Synthesis runs on a worker thread because \texttt{\_speak} shares an event loop with \texttt{ingest\_frame}: synthesis is sub-second but still blocking, and perception should not stall mid-mission for a voice line. A failure degrades the mission to text-only rather than stopping it, the same soft-fail stance the pre-departure dialogue's own \texttt{\_speak} in \texttt{pipeline.py} takes.}
  \label{fig:flow-tts}
\end{figure}

%% ===================================================================
\mysection{Model inference layer}{Model Inference Layer}
\label{sec:inference}

\mysubsection{Planner}{Planner}
\label{sec:planner}

\texttt{InferenceProvider} is an abstract class with one method,
\texttt{complete}, that takes a list of chat messages and an optional image
and returns a \texttt{Completion}: the text, which provider and model
answered, the latency, and token counts and cost when the provider reports
them. \texttt{OllamaProvider} and \texttt{OpenAICompatibleProvider} both
implement it, the second covering OpenAI and DeepSeek, since both speak the
same \texttt{/chat/completions} schema; a DeepSeek instance is constructed
with \texttt{supports\_vision=False}, and \texttt{get\_provider} refuses to
hand a text-only provider to the VLM role. Every construction is cached by
(role, provider, model), so a long-running process reuses one client rather
than reconnecting on every call, and swapping the planner or the vision
model is a matter of an environment variable
(\texttt{PLANNER\_PROVIDER}, \texttt{VLM\_PROVIDER}), not a code change.

\begin{figure}[htbp]
  \centering
  \begin{tikzpicture}[node distance=0.8cm]
    \node[blk] (pl) at (0,1.2)    {\texttt{planner.py}};
    \node[blk] (vl) at (0,-1.2)   {\texttt{vlm.py}};
    \node[mdl] (if) at (4.0,0)    {\texttt{InferenceProvider}\\\texttt{.complete()}};
    \node[mdl] (ol) at (8.5,1.2)  {\texttt{OllamaProvider}\\(local)};
    \node[mdl] (oa) at (8.5,-1.2) {\texttt{OpenAICompatibleProvider}\\(OpenAI, DeepSeek)};
    \draw[arr] (pl)--(if);
    \draw[arr] (vl)--(if);
    \draw[arr] (if)--(ol);
    \draw[arr] (if)--(oa);
  \end{tikzpicture}
  \caption{Both roles call the same interface. \texttt{get\_provider(role)} resolves \texttt{PLANNER\_PROVIDER} or \texttt{VLM\_PROVIDER} to a concrete class and caches the instance.}
  \label{fig:providers}
\end{figure}

\texttt{OllamaProvider} sends \texttt{num\_ctx} (8192 by default), an
optional \texttt{repeat\_penalty}, and \texttt{format="json"} for every
planner and VLM call. Two failures are retried once before being allowed to
fail a mission: a runner crash during a cold CUDA load on
\texttt{gemma4:e4b}, which happened during testing on 5 October, and
Ollama's own repeat-limit abort, which happens when a small model loops on
the same tokens. \texttt{build\_plan\_prompt} assembles the prompt from the
active planner profile (Section~\ref{sec:vocabulary}), the room digest when
one is loaded, and the conditional guidance block when the command contains
a condition word (Section~\ref{sec:conditional}); it is a pure function,
separate from \texttt{generate\_plan} itself, so a small tool,
\texttt{tools/snapshot\_prompt.py}, can print the exact prompt a command
would produce without calling a model.

\begin{figure}[htbp]
  \centering
  \begin{tikzpicture}
    \node[term] (a) at (0,0)     {command, profile,\\room digest};
    \node[proc] (b) at (0,-0.9)  {\texttt{build\_plan\_prompt()}};
    \node[proc] (c) at (0,-1.8)  {\texttt{provider.complete()}};
    \node[dec]  (d) at (-3.4,-3.3) {step is a\\well-formed\\object?};
    \node[proc] (e) at (0,-3.3)  {parse JSON reply};
    \node[proc] (f) at (-3.4,-4.8) {drop step;\\note in plan};
    \node[proc] (g) at (0,-5.4)  {\texttt{\_promote\_moves\_}\\\texttt{to\_approach}};
    \node[proc] (h) at (0,-6.3)  {\texttt{validate\_steps}\\(role-specific actions)};
    \node[term] (i) at (0,-7.2)  {return plan};
    \draw[fa] (a)--(b); \draw[fa] (b)--(c); \draw[fa] (c)--(e); \draw[fa] (e)--(d);
    \draw[fa] (d.west) -- node[fl,above]{no} (f.east);
    \draw[fa] (d.south) -- node[fl,right]{yes} (g.north);
    \draw[fa] (f) -- (g.west);
    \draw[fa] (g)--(h); \draw[fa] (h)--(i);
  \end{tikzpicture}
  \caption{\texttt{generate\_plan}. A smaller model sometimes writes a step as a bare string instead of an object, or nests the list oddly, so every step is checked before anything downstream calls \texttt{.get} on it. \texttt{\_promote\_moves\_to\_approach} rewrites a \texttt{move} step whose target names an object into an \texttt{approach} step (Section~\ref{sec:vocabulary}); the retried transient failures from \texttt{provider.complete()} are omitted here for space and described in the text above.}
  \label{fig:flow-planner}
\end{figure}

Failure recovery (Section~\ref{ssec:guidance_loop}) is a second prompt
builder, \texttt{build\_replan\_prompt}, and a second entry point,
\texttt{replan\_from\_state}, rather than a parameter on
\texttt{generate\_plan}, because a replan is built from the command, the
steps already done, the step that failed and why, and the guidance
conversation so far, which a first plan has no use for.

\mysubsection{Vision-language verification}{Vision-language Verification}
\label{sec:vlm}

Every call into the vision model goes through one function, \texttt{\_ask},
which takes a named stage (for logging), a prompt, an optional image and a
salvage flag, resolves the VLM provider, and sends the prompt with
\texttt{json\_mode=True}. Every stage built on it is a separate, short,
single-purpose prompt: \texttt{check\_ambiguity}, \texttt{verify\_plan},
\texttt{describe\_frame}, \texttt{inventory\_frame}, \texttt{locate\_target},
\texttt{check\_condition} and \texttt{survey\_frame}. None of them share a
prompt template, deliberately, since a prompt that tries to serve several
callers tends to grow extra fields that most callers never fill in, which
makes a wrong answer harder to trace back to a cause.

\begin{figure}[htbp]
  \centering
  \begin{tikzpicture}
    \node[term] (a) at (0,0)    {stage, prompt,\\image, salvage};
    \node[proc] (b) at (0,-0.9) {\texttt{provider.complete()}\\(\texttt{json\_mode=True})};
    \node[dec]  (c) at (0,-2.3) {parses as\\JSON?};
    \node[term] (d) at (2.9,-2.3) {return\\parsed reply};
    \node[dec]  (e) at (0,-3.9) {\texttt{salvage=True}?};
    \node[term] (f) at (2.9,-3.9) {raise;\\stage fails};
    \node[proc] (g) at (0,-5.3) {\texttt{\_salvage\_truncated\_json}:\\try last 20 cut points};
    \node[dec]  (h) at (0,-6.9) {a cut point\\parses?};
    \node[term] (i) at (-2.9,-8.2) {return\\partial result};
    \node[term] (j) at (2.9,-8.2) {raise;\\stage fails};
    \draw[fa] (a)--(b); \draw[fa] (b)--(c);
    \draw[fa] (c.east) -- node[fl,above]{yes} (d.west);
    \draw[fa] (c.south) -- node[fl,left]{no} (e.north);
    \draw[fa] (e.east) -- node[fl,above]{no} (f.west);
    \draw[fa] (e.south) -- node[fl,left]{yes} (g.north);
    \draw[fa] (g)--(h);
    \draw[fa] (h.south west) -- node[fl,above]{yes} (i.north);
    \draw[fa] (h.south east) -- node[fl,above]{no} (j.north);
  \end{tikzpicture}
  \caption{\texttt{\_ask}. Salvage is only enabled for \texttt{inventory\_frame} and \texttt{survey\_frame}, the two calls that ask for a list of objects and sometimes hit \texttt{num\_predict} mid-list on a model that has started repeating itself; a reply with ten correct objects and one truncated one is more useful than discarding it. Every other call keeps \texttt{salvage=False}, since a truncated yes/no answer or bounding box is not something a partial parse can safely recover.}
  \label{fig:flow-vlm}
\end{figure}

\texttt{locate\_target} is the one with the most moving parts, since it is
called once per cycle of the free-roam approach (Section~\ref{ssec:freeroam}).
It asks for a bounding box, a confidence label, and a distance estimate;
\texttt{distance\_cm\_from} converts whichever unit the model used into
centimetres, or returns \texttt{None} for any value outside
\((0,\;\texttt{VLM\_DIST\_MAX\_CM}]\), including a \texttt{NaN} that would
otherwise fail every comparison silently. Every numeric field from the model
is treated this way, accepted only inside a sane range, because a model this
size occasionally returns a plausible-looking number for a question it
cannot actually answer.

%% ===================================================================
\mysection{Clarification and mission orchestration}{Clarification and Mission Orchestration}
\label{sec:orchestration}

\mysubsection{Clarification manager}{Clarification Manager}
\label{sec:clarification_impl}

The clarification loop (Section~\ref{sec:clarification_design}) is a small
state machine, \texttt{DialogueSession}, with \texttt{Phase} as a plain
enumeration (\texttt{CLARIFYING}, \texttt{PLANNING}, \texttt{VERIFYING},
\texttt{AWAITING\_CONFIRMATION}, \texttt{EXECUTING}, \texttt{REPORTING},
\texttt{CANCELLED}), and a module of handler functions in
\texttt{pipeline.py} that mutate it. \texttt{effective\_command} is the one
property everything downstream reads instead of the raw transcript: the
resolved command once the vision model has rewritten one, the original
command once the session is capped, and otherwise the command with the
dialogue so far appended, in that order. Confirmation is a smaller state
machine of its own, in \texttt{confirmation.py}: \texttt{keyword\_pass}
checks the reply against an English and an Afrikaans phrase list, reject
checked first, and only falls back to a planner-model call,
\texttt{classify}, when neither matches.

\begin{figure}[htbp]
  \centering
  \begin{tikzpicture}
    \node[term] (a) at (0,0)    {\texttt{effective\_command}};
    \node[proc] (b) at (0,-0.9) {\texttt{check\_ambiguity()}};
    \node[dec]  (c) at (0,-2.3) {still\\ambiguous?};
    \node[term] (d) at (-3.0,-3.9) {move to\\confirmation};
    \node[dec]  (e) at (3.6,-2.3) {asked\\already?};
    \node[term] (f) at (3.6,-3.9) {mark\\\texttt{capped}; stop};
    \node[term] (g) at (7.6,-2.3) {append question;\\ask it};
    \draw[fa] (a)--(b); \draw[fa] (b)--(c);
    \draw[fa] (c.south west) -- node[fl,above]{no} (d.north);
    \draw[fa] (c.east) -- node[fl,above]{yes} (e.west);
    \draw[fa] (e.south) -- node[fl,left]{yes} (f.north);
    \draw[fa] (e.east) -- node[fl,above]{no} (g.west);
  \end{tikzpicture}
  \caption{\texttt{\_advance\_to\_confirmation}, one turn of the clarification loop. The repeat-question check exists because a model that cannot think of anything new to ask would otherwise repeat the same question until the five-turn cap was reached, rather than giving up as soon as it had nothing left to add.}
  \label{fig:flow-clarify}
\end{figure}

\texttt{DialogueSession} hands off to \texttt{MissionSession} once
confirmation succeeds, and everything from there on, including recovery and
revision, is the orchestration described in the rest of this section.

\mysubsection{The mission execution loop}{The Mission Execution Loop}
\label{ssec:mission_loop}

Once a plan is confirmed, \texttt{run\_mission} is the single function that
drives it, one step at a time, until the mission reaches a terminal phase or
runs out of steps. Which function actually runs a step depends on what it is
and what the rover can do: a \texttt{check} step goes to
\texttt{\_run\_check}; an \texttt{approach} step goes to
\texttt{approach\_controller} (the free-roam hop loop,
Section~\ref{ssec:freeroam}), but only when \texttt{NAV\_MODE} is
\texttt{free} and the rover object exposes \texttt{pivot}, \texttt{hop},
\texttt{get\_frame} and \texttt{set\_gimbal}; everything else goes straight
to \texttt{rover.execute\_step} on a worker thread.

\begin{figure}[htbp]
  \centering
  \begin{tikzpicture}
    \node[term] (a) at (0,0)     {next step\\(by cursor)};
    \node[dec]  (b) at (0,-1.4)  {conditional tag\\matches recorded\\branch?};
    \node[proc] (c) at (-3.4,-1.4) {skip step;\\cursor++};
    \node[proc] (d) at (0,-3.0)  {dispatch: \texttt{check} /\\free-approach /\\\texttt{execute\_step}};
    \node[dec]  (e) at (0,-4.6)  {result\\status?};
    \node[proc] (f) at (-3.6,-6.0) {\texttt{ok}: post-\\process; cursor++};
    \node[proc] (g) at (0,-6.0)  {\texttt{blocked}:\\guidance loop\\(Fig.~\ref{fig:flow-guidance})};
    \node[dec]  (h) at (0,-7.8)  {replanned?};
    \node[term] (i) at (3.6,-6.0)  {\texttt{halted}:\\mission ends};
    \draw[fa] (a)--(b);
    \draw[fa] (b.west) -- node[fl,above]{no} (c.east);
    \draw[fa] (c) |- (a);
    \draw[fa] (b.south) -- node[fl,left]{yes} (d.north);
    \draw[fa] (d)--(e);
    \draw[fa] (e.south west) -- node[fl,above]{ok} (f.north);
    \draw[fa] (e.south) -- node[fl,right]{blocked} (g.north);
    \draw[fa] (e.south east) -- node[fl,above]{halted} (i.north);
    \draw[fa] (g)--(h);
    \draw[fa] (f) |- (a);
    \draw[fa] (h.west) -- node[fl,above]{yes} ++(-5.2,0) |- (a);
    \draw[fa] (h.south) -- node[fl,right]{no} ++(0,-1.0) node[term]{mission ends};
  \end{tikzpicture}
  \caption{\texttt{run\_mission}'s step loop. The conditional-tag check is what Section~\ref{sec:conditional} relies on: the planner only decides which steps carry a \texttt{when} tag, and this check is what honours it. The loop's own \texttt{finally} block (not shown) writes the report for every non-aborted outcome; an operator-triggered \texttt{asyncio.CancelledError} skips straight to \texttt{rover.halt()} instead.}
  \label{fig:flow-mission}
\end{figure}

A \texttt{blocked} result hands off to the guidance loop
(Section~\ref{ssec:guidance_loop}): a staged replan resumes the loop from the
top with the cursor reset to~0; otherwise the mission ends. A \texttt{halted}
result ends the mission immediately with no guidance offered, since that
status means the step itself decided the mission cannot continue, a
narrower condition than being blocked and askable. Every non-aborted mission
writes its report in the loop's \texttt{finally} block, whether it finished
cleanly, ran out of guidance, or was halted outright, because a run that
ended badly is exactly the run whose report matters most for the
evaluation. The one path that skips this is an operator-triggered
\texttt{asyncio.CancelledError}: that path calls \texttt{rover.halt()}
immediately and writes the deterministic half of the report straight to
disk with \texttt{spoken=False}, so stopping the rover is never delayed by a
planner call or a spoken summary it has no time to give.

\mysubsection{Failure recovery in code}{Failure Recovery in Code}
\label{ssec:guidance_loop}

A \texttt{blocked} result is handed to \texttt{\_guidance\_loop}, which is
where the recovery behaviour of Section~\ref{sec:recovery} is actually
implemented. If \texttt{GUIDANCE\_ENABLED} is off, it halts immediately,
mainly for test batches that want a plain pass/fail rather than a multi-turn
conversation. Otherwise it builds a question from a fixed template keyed on
the failure reason (\texttt{target\_not\_found}, \texttt{obstacle},
\texttt{arrival\_contradicted}, \texttt{cycle\_budget\_exhausted}, or a
generic fallback) and asks it, up to \texttt{MAX\_GUIDANCE\_TURNS} times
(3 by default).

\begin{figure}[htbp]
  \centering
  \begin{tikzpicture}
    \node[term] (a) at (0,0)    {step blocked};
    \node[dec]  (b) at (0,-1.4) {\texttt{GUIDANCE\_}\\\texttt{ENABLED}?};
    \node[term] (c) at (3.2,-1.4) {halt;\\mission ends};
    \node[proc] (d) at (0,-3.0) {ask question;\\wait for answer};
    \node[dec]  (e) at (0,-4.6) {declines or\\timed out?};
    \node[term] (f) at (3.2,-4.6) {halt with reason\\(timeout/declined);\\mission ends};
    \node[proc] (g) at (0,-6.2) {\texttt{replan\_from\_state}};
    \node[dec]  (h) at (0,-7.8) {usable steps\\returned?};
    \node[term] (i) at (-3.2,-9.2) {stage\\\texttt{pending\_swap};\\return \texttt{True}};
    \node[term] (j) at (3.2,-9.2) {3 turns\\used up?};
    \draw[fa] (a)--(b);
    \draw[fa] (b.east) -- node[fl,above]{no} (c.west);
    \draw[fa] (b.south) -- node[fl,left]{yes} (d.north);
    \draw[fa] (d)--(e);
    \draw[fa] (e.east) -- node[fl,above]{yes} (f.west);
    \draw[fa] (e.south) -- node[fl,left]{no} (g.north);
    \draw[fa] (g)--(h);
    \draw[fa] (h.south west) -- node[fl,above]{yes} (i.north);
    \draw[fa] (h.east) -- node[fl,above]{no} (j.west);
    \draw[fa] (j.north) |- ++(0,0.6) -| (d.east);
  \end{tikzpicture}
  \caption{\texttt{\_guidance\_loop}. A ``no'' out of the turn-count check on the right loops back to ask again rather than ending the mission; running out of turns is itself a third halt reason (\texttt{guidance\_exhausted}), alongside timeout and decline, so a batch of failed runs can be told apart by why they failed.}
  \label{fig:flow-guidance}
\end{figure}

An answer that declines is checked with \texttt{\_declines}, not
\texttt{confirmation.keyword\_pass}, because the confirmation reject list
includes ordinary words like ``no'' and ``wrong'' that show up constantly
inside a perfectly good guidance answer such as ``there is no box on the
left, try the shelf''. \texttt{\_declines} only triggers on an answer that
opens with an explicit stop word, or a bare answer of three words or fewer
that \texttt{confirmation.keyword\_pass} alone calls a reject. In tests,
\texttt{mission.guidance\_provider} can be set to a plain synchronous
function instead of waiting on audio or a WebSocket reply, which is how a
batch of L0 to L5 runs scripts an answer for every question the loop asks.

\mysubsection{Mid-mission revision}{Mid-mission Revision}
\label{ssec:revision_impl}

Every frame that arrives during a mission passes through
\texttt{ingest\_frame}, which gates three tiers of increasingly expensive
work so the VLM is not asked to look at every frame a camera happens to
send.

\begin{figure}[htbp]
  \centering
  \begin{tikzpicture}
    \node[term] (a) at (0,0)    {frame arrives};
    \node[dec]  (b) at (0,-1.4) {tier 0: sharp\\and changed\\enough?};
    \node[term] (c) at (3.0,-1.4) {drop frame};
    \node[proc] (d) at (0,-3.0) {tier 1: caption\\(always recorded)};
    \node[dec]  (e) at (0,-4.4) {mission is\\\texttt{EXECUTING}?};
    \node[term] (f) at (3.0,-4.4) {stop here};
    \node[proc] (g) at (0,-6.0) {tier 2: \texttt{revise\_plan}};
    \node[dec]  (h) at (0,-7.4) {change\\proposed?};
    \node[term] (i) at (3.0,-7.4) {no action};
    \node[proc] (j) at (0,-9.0) {\texttt{revision.classify}:\\NO\_CHANGE / REROUTE /\\MATERIAL / BLOCKED};
    \node[term] (k) at (0,-10.4) {apply, confirm,\\or halt accordingly};
    \draw[fa] (a)--(b);
    \draw[fa] (b.east) -- node[fl,above]{no} (c.west);
    \draw[fa] (b.south) -- node[fl,left]{yes} (d.north);
    \draw[fa] (d)--(e);
    \draw[fa] (e.east) -- node[fl,above]{no} (f.west);
    \draw[fa] (e.south) -- node[fl,left]{yes} (g.north);
    \draw[fa] (g)--(h);
    \draw[fa] (h.east) -- node[fl,above]{no} (i.west);
    \draw[fa] (h.south) -- node[fl,left]{yes} (j.north);
    \draw[fa] (j)--(k);
  \end{tikzpicture}
  \caption{\texttt{ingest\_frame}'s three tiers, feeding \texttt{revision.classify}. Tier 0 is a 32$\times$32 greyscale thumbnail diff plus a Laplacian-variance blur score (\texttt{PERCEPTION\_MIN\_SHARPNESS}, \texttt{PERCEPTION\_MIN\_INTERVAL\_S}, \texttt{PERCEPTION\_CHANGE\_THRESHOLD}); tier 1 is one \texttt{describe\_frame} caption; tier 2 only fires while executing.}
  \label{fig:flow-revision}
\end{figure}

\texttt{revision.classify} is a plain rule, not a model call: it compares
the confirmed plan's targets against the proposed plan's and the action
vocabulary each uses. Adding a target, dropping a confirmed one, or
introducing an action outside a fixed navigation whitelist is
\texttt{MATERIAL} and stops the mission for a spoken yes before
\texttt{handle\_revision\_confirmation} applies it; an empty revised plan is
\texttt{BLOCKED} and ends the mission outright; only when the goal and every
target are unchanged, and the only new actions are navigation ones, is the
change a silent \texttt{REROUTE}. I kept this classification out of the
model deliberately, for the same reason \texttt{confirmation.keyword\_pass}
exists ahead of a model fallback: whether a human must be consulted is not a
decision to delegate to something that can hallucinate its way to ``no
change needed''.

%% ===================================================================
\mysection{Rover client and reactive layer}{Rover Client and Reactive Layer}
\label{sec:rover_client}

\mysubsection{Real hardware: PiRoverController}{Real Hardware}
\label{ssec:pi_rover}

\texttt{RoverController} is an abstract class with two required methods,
\texttt{execute\_step} and \texttt{halt}, and one optional one,
\texttt{status\_light}, defaulting to a no-op. Three classes satisfy it:
\texttt{SimulatedRover}, a timer that always reports success;
\texttt{VirtualRover}, described in Section~\ref{ssec:virtual_rover}; and
\texttt{PiRoverController}, which drives the real TurboPi over
\texttt{roslibpy}. On construction it opens one \texttt{roslibpy.Ros}
connection, advertises a \texttt{/cmd\_vel} topic for driving and a
PWM-servo topic for the gimbal, and opens two ROS2 services,
\texttt{/line\_follow/cmd} and \texttt{/line\_follow/set\_running}, that
control the stock line-follower node. Two threading events carry meanings
that are easy to conflate: \texttt{\_halted} means the controller has been
told to stop and is cleared only by the next \texttt{execute\_step}, while
\texttt{\_target\_confirmed} means a mid-drive scan has confirmed the named
target is visible, a different event from being stopped.

\begin{figure}[htbp]
  \centering
  \begin{tikzpicture}
    \node[term] (a) at (0,0)    {drive forward\\(stock line node)};
    \node[dec]  (b) at (0,-1.4) {services\\reachable?};
    \node[term] (c) at (3.3,-1.4) {wait 0.3\,s;\\\texttt{blocked}:\\\texttt{line\_follow\_}\\\texttt{unavailable}};
    \node[dec]  (d) at (0,-3.2) {arrived\\(\texttt{crossroads\_stop}\\or confirmed)?};
    \node[term] (e) at (3.3,-3.2) {\texttt{ok}: arrived};
    \node[dec]  (f) at (0,-5.0) {sonar below\\threshold?};
    \node[term] (g) at (3.3,-5.0) {\texttt{blocked}:\\obstacle};
    \node[dec]  (h) at (0,-6.8) {timeout\\exceeded?};
    \node[term] (i) at (3.3,-6.8) {\texttt{blocked}:\\timeout};
    \draw[fa] (a)--(b);
    \draw[fa] (b.east) -- node[fl,above]{no} (c.west);
    \draw[fa] (b.south) -- node[fl,left]{yes} (d.north);
    \draw[fa] (d.east) -- node[fl,above]{yes} (e.west);
    \draw[fa] (d.south) -- node[fl,left]{no} (f.north);
    \draw[fa] (f.east) -- node[fl,above]{yes} (g.west);
    \draw[fa] (f.south) -- node[fl,left]{no} (h.north);
    \draw[fa] (h.east) -- node[fl,above]{yes} (i.west);
    \draw[fa] (h.south) -- node[fl,left]{no} ++(0,-0.6) -- ++(-2.4,0) |- (d.west);
  \end{tikzpicture}
  \caption{\texttt{\_follow\_line}, the supervisor added around the stock line-follower node. It owns exactly the three things the node has no opinion about: arrival, obstacle stopping through the polled sonar (using \texttt{\_sonar\_seq} to tell two fresh readings from one stale one read twice), and a wall-clock give-up, since the node itself drives forward indefinitely if the line is lost. The \texttt{\_lf\_unavailable} check (added 11 October) means a drive that never started is reported as blocked rather than as a drive that ran and found nothing.}
  \label{fig:flow-followline}
\end{figure}

\mysubsection{Virtual testing: VirtualRover}{Virtual Testing}
\label{ssec:virtual_rover}

\texttt{VirtualRover} is what every run without physical hardware actually
executes against, built around one idea: \texttt{approach} is a navigation
skill, so it plans a route and follows it, while \texttt{move} and
\texttt{turn} are literal, executed exactly as commanded with no obstacle
avoidance of their own. Both halves read from the same
\texttt{OccupancyGrid} (\texttt{nav.py}): every obstacle polygon plus the
room boundary, inflated by the rover's footprint clearance (computed once
from its length and width, 18.7\,cm by 16.2\,cm by default, plus a safety
margin), so a planner and a collision checker built from two different
descriptions of the room can never disagree about what is free.

\begin{figure}[htbp]
  \centering
  \begin{tikzpicture}
    \node[dec]  (a) at (0,0)      {step\\action?};
    \node[proc] (b) at (-5.2,-1.6) {\texttt{approach}: resolve\\target, A* plan\\(\texttt{nav.plan\_path}),\\walk waypoints};
    \node[proc] (c) at (-1.7,-1.6) {\texttt{move}/\texttt{turn}:\\magnitude + drift,\\\texttt{nav.sweep\_segment}};
    \node[proc] (d) at (1.7,-1.6) {\texttt{scan}/\texttt{stop}/\\\texttt{report}:\\fixed dwell};
    \node[proc] (e) at (5.2,-1.6)  {unknown\\action};
    \node[term] (f) at (0,-3.2)   {return \texttt{ok} /\\\texttt{blocked} /\\\texttt{halted}};
    \draw[fa] (a) -- node[fl,above,sloped]{approach} (b.north);
    \draw[fa] (a) -- node[fl,above,sloped]{move/turn} (c.north);
    \draw[fa] (a) -- node[fl,above,sloped]{scan etc.} (d.north);
    \draw[fa] (a) -- node[fl,above,sloped]{else} (e.north);
    \draw[fa] (b)--(f); \draw[fa] (c)--(f); \draw[fa] (d)--(f); \draw[fa] (e)--(f);
  \end{tikzpicture}
  \caption{\texttt{VirtualRover.execute\_step}'s dispatch. An \texttt{approach} step can only fail by planning (no path, or an unresolved target logged as a hallucination for RQ2), never by collision; a \texttt{move} or \texttt{turn} step is swept against the same grid and stops at the last free point if blocked, so a collision there means the chain of plan steps itself went somewhere it should not have.}
  \label{fig:flow-virtualrover}
\end{figure}

The room itself, a \texttt{RoomMap} loaded from JSON
(\texttt{room\_map.py}), is hand-typed, not estimated: coordinates for every
named landmark are measured off a sketch and a tape measure once per test
room, with a landmark's point being the standing spot a rover can occupy to
look at the thing, not the thing's own centroid. Every step also applies a
configurable drift, Gaussian jitter on the achieved position and heading,
scaled by \texttt{VIRTUAL\_DRIFT\_FRAC} and \texttt{VIRTUAL\_DRIFT\_DEG}; both
are left at zero in the evaluation sweeps, so the numbers in
Chapter~\ref{chap:evaluation} are not inflated or deflated by injected
noise. Simulated and wall-clock time are kept apart the same way: every
action's duration is computed in simulated time first and only then divided
by \texttt{VIRTUAL\_TIME\_SCALE}, so a sweep of many runs can be played back
faster than real time without changing the distances, speeds, or timings
the evaluation reports on. Every run writes one JSON line per step to a
trace file, and \texttt{summary()} rolls a run's trace up into the numbers
RQ2 is built on: step count, how many were blocked, how many targets were
hallucinated and which ones, collisions, how often no path existed at all,
and the run's final pose.

%% ===================================================================
\mysection{Communication protocol and dashboard}{Communication Protocol and Dashboard}
\label{sec:protocol}

The brain is a FastAPI application. The two endpoints that matter for a
mission, \texttt{/ws/dialogue} and \texttt{/ws/execution}, are WebSockets
carrying small JSON frames rather than a request-response API, since both
sides need to push events the other did not ask for. This is a message
protocol rather than a branching algorithm, so Figure~\ref{fig:protocol}
below is a sequence diagram, not a flowchart: it shows who sends what, not a
decision to be made.

\begin{figure}[htbp]
  \centering
  \begin{tikzpicture}[node distance=1.0cm]
    \node[hw] (pi) at (0,0)   {Pi client};
    \node[blk] (br) at (8,0)  {Brain (FastAPI)};
    \draw[arr] ([yshift=0.9cm]pi.east) -- node[lbl, above]{start, audio\_chunk, confirmation\_audio}([yshift=0.9cm]br.west);
    \draw[arr] ([yshift=0.3cm]br.west) -- node[lbl, above]{session, speak, plan\_ready, execute, cancelled}([yshift=0.3cm]pi.east);
    \draw[arr] ([yshift=-0.3cm]pi.east) -- node[lbl, below]{begin, frame, revision\_audio}([yshift=-0.3cm]br.west);
    \draw[arr] ([yshift=-0.9cm]br.west) -- node[lbl, below]{step\_started, check, awaiting\_guidance, report}([yshift=-0.9cm]pi.east);
  \end{tikzpicture}
  \caption{\texttt{/ws/dialogue} (top pair) carries everything up to confirmation; \texttt{/ws/execution} (bottom pair) carries the mission. One frame per complete utterance, not a stream.}
  \label{fig:protocol}
\end{figure}

Both connections are opened with \texttt{max\_size=None}. This was not the
original setting; a long spoken plan read-back, a \texttt{speak} event
carrying the base64 WAV Section~\ref{ssec:tts} describes, exceeded the
websockets library's 1\,MiB default and closed the connection with code
1009, which I only found by reproducing a failure that happened on longer
plans and not shorter ones. \texttt{/scene} accepts an uploaded walkthrough
video and returns a scene id; \texttt{/logs} and
\texttt{/logs/\{session\_id\}} serve the run log listing and detail the
dashboard's Logs page reads; and \texttt{/logs/\{session\_id\}/flag} is a
small \texttt{PUT} endpoint that lets an operator mark a run good, bad or
for review with a note, stored separately in \texttt{run\_flags.json} so
flagging a run never touches the record the run itself wrote.

The dashboard is a React and TypeScript single-page app, a thin client over
these same endpoints, holding no logic of its own that the evaluation
depends on: every number Chapter~\ref{chap:evaluation} reports comes from
the JSON logs, not from the UI. It is not part of the project's
contribution, so I leaned on an AI assistant to build most of its
components rather than writing them by hand. It is a debugging and
demonstration tool, nothing more: something to let me upload a room
walkthrough video, watch a dialogue and a mission happen panel by panel,
and drive the rover manually, without needing a terminal and a stack of
\texttt{curl} commands to check that the pipeline underneath was actually
working end to end.

%% ===================================================================
\mysection{Logging and measurement instrumentation}{Logging and Measurement Instrumentation}
\label{sec:logging}

Two logs exist side by side and serve different purposes. \texttt{brain.log}
is a conventional application log, written through Python's \texttt{logging}
module, interleaving every call in time order with full prompts and raw
model output at debug level; it is for debugging a run, not for the
evaluation, and I do not parse it programmatically. The run log, one JSON
file per mission, is the one the evaluation is built on, and
\texttt{report.py} is what assembles it: the plan, every check and its
result, the guidance log, the light-state timeline, the model tag for both
roles, and a tagged \texttt{test\_level} when one is set, so a batch of L0
to L5 runs can be grouped without reading every command by hand.

\begin{figure}[htbp]
  \centering
  \begin{tikzpicture}
    \node[term] (a) at (0,0)    {mission ends};
    \node[dec]  (b) at (0,-1.4) {guidance turns $>$ 0\\and phase \texttt{completed}\\or \texttt{halted}?};
    \node[term] (c) at (-3.2,-2.9) {outcome =\\phase};
    \node[term] (d) at (3.2,-2.9) {outcome = phase +\\\texttt{\_after\_guidance}};
    \node[dec]  (e) at (3.2,-4.6) {\texttt{approach}\\succeeded after\\last \texttt{blocked}\\step?};
    \node[term] (f) at (6.4,-4.6) {\texttt{recovered}\\= True};
    \node[term] (g) at (3.2,-6.2) {\texttt{recovered}\\= False};
    \draw[fa] (a)--(b);
    \draw[fa] (b.south west) -- node[fl,above]{no} (c.north);
    \draw[fa] (b.south east) -- node[fl,above]{yes} (d.north);
    \draw[fa] (d)--(e);
    \draw[fa] (e.east) -- node[fl,above]{yes} (f.west);
    \draw[fa] (e.south) -- node[fl,left]{no} (g.north);
  \end{tikzpicture}
  \caption{\texttt{report.py}'s outcome classification. A mission that needed at least one guidance turn is never reported as a plain \texttt{completed}, even if it finished, because the distinction between a clean run and a recovered one (Section~\ref{sec:recovery}) is exactly what RQ1 needs to see; \texttt{recovered} only counts a run that actually reached its target afterwards, not merely one that ended without halting.}
  \label{fig:flow-outcome}
\end{figure}

Every model call also writes one entry to a smaller usage ledger
(\texttt{providers/usage.py}), holding latency, token counts and cost when
the provider reports them, reset at the start of every session. A separate,
offline tool, \texttt{tools/context\_budget.py}, measures something the run
log cannot: how many tokens a given prompt would cost before it is ever
sent. It imports \texttt{planner.py} and \texttt{vlm.py} directly, forces
\texttt{ROVER=virtual}, and asks Ollama for the \texttt{prompt\_eval\_count}
on a constructed prompt with \texttt{num\_predict=1}, which counts the
tokens without paying for a real generation, measuring the fixed
instructions, the room digest, one clarification turn and one recovery turn
separately.
